% GENERATED FILE. Edit canonical lessons, then run npm run generate:books.
% canonical-source-hash: fnv1a64:19d8ad1e6c9fb9ef
% canonical-lessons: TA-C533-pantayam, TA-C533-turattu, TA-C533-carukku, TA-C533-tavu, TA-C533-mukkuli

\chapter{The Race}
\label{ch:ta-the-race}
\hlchaptermodality{533}

\begin{chapteropening}
\textbf{By the end of this chapter:} \emph{I can talk about racing, chasing and diving.}

Complete the last lesson of chapter 533: I can talk about racing, chasing and diving.
\end{chapteropening}

\section[\texorpdfstring{\ta{பந்தயம்} (pantayam)}{பந்தயம் (pantayam)}]{\ta{பந்தயம்} (pantayam) — a race; a bet}

\label{lesson:TA-C533-pantayam}



\begin{quote}
Before the new one: say the Tamil for to kick, then the Tamil for to hurl.
\end{quote}

\subsection*{You'll want to know: \ta{பந்தயம்}}
\textbf{\ta{பந்தயம்}} — \emph{pantayam} — \textquotedblleft{}a race; a bet\textquotedblright{}.

\begin{quote}
\ta{நூறு} \ta{மீட்டர்} \ta{பந்தயத்தில்} \ta{அவன்} \ta{முதலில்} \ta{வந்தான்}.
\end{quote}

\begin{quote}
\emph{nūṟu mīṭṭar pantayattil avaṉ mutalil vantāṉ.}
\end{quote}

\begin{quote}
He came first in the hundred-metre race.
\end{quote}

Both \textquotedblleft{}a race\textquotedblright{} and \textquotedblleft{}a bet\textquotedblright{}: \ta{பந்தயம்} \ta{கட்டு}, to place a bet.

\begin{cousinweb}
DEDR 3921 gives \textbf{\ta{பந்தயம்}} (\emph{pantayam}), \textquotedblleft{}a contest for a prize, a wager\textquotedblright{}, with cognates across the south, and reports a proposal to derive the group from Sanskrit \emph{paṇita}, \textquotedblleft{}a wager\textquotedblright{}; the origin is not settled.
\end{cousinweb}

\begin{grammarlens}[title={What you've built}]
One more word for describing what you lived and felt.
\end{grammarlens}

\subsection*{Your turn}
\begin{itemize}
\raggedright
  \item \emph{Say it:} \emph{pantayam}
  \item \emph{Say it:} \emph{pantayam}, once more
  \item \emph{Say it:} \emph{nūṟu mīṭṭar pantayattil avaṉ mutalil vantāṉ.} — He came first in the hundred-metre race.
  \item \emph{Recall:} say the Tamil for to kick, then the Tamil for to hurl, then say \emph{pantayam} again
\end{itemize}

\begin{tcolorbox}[breakable,colback=teal!4,colframe=teal!35!black,title={Before you move on}]
What does \ta{பந்தயம்} mean? (\textquotedblleft{}A race; a bet\textquotedblright{}.) Say it once more.
\end{tcolorbox}
\section[\texorpdfstring{\ta{துரத்து} (turattu)}{துரத்து (turattu)}]{\ta{துரத்து} (turattu) — to chase; to drive away}

\label{lesson:TA-C533-turattu}



\begin{quote}
Before the new one: say the Tamil for to hurl, then the Tamil for a race.
\end{quote}

\subsection*{You'll want to know: \ta{துரத்து}}
\textbf{\ta{துரத்து}} — \emph{turattu} — \textquotedblleft{}to chase; to drive away\textquotedblright{}.

\begin{quote}
\ta{நாய்} \ta{பூனையைத்} \ta{துரத்தியது}.
\end{quote}

\begin{quote}
\emph{nāy pūṉaiyait turattiyatu.}
\end{quote}

\begin{quote}
The dog chased the cat.
\end{quote}

Past (I …): \ta{துரத்தினேன்}, \emph{turattiṉēṉ}.

What is chased takes -\ta{ஐ}: \ta{பூனையைத்} \ta{துரத்து}.

\begin{cousinweb}
\textbf{\ta{துரத்து}} (\emph{turattu}) is native Dravidian: DEDR 3340 gives \ta{துர} (\emph{tura}), \textquotedblleft{}to drive, to drive away\textquotedblright{}, and \ta{துரத்து} (\emph{turattu}), \textquotedblleft{}to chase\textquotedblright{}.
\end{cousinweb}

\begin{grammarlens}[title={What you've built}]
One more word for describing what you lived and felt.
\end{grammarlens}

\subsection*{Your turn}
\begin{itemize}
\raggedright
  \item \emph{Say it:} \emph{turattu}
  \item \emph{Say it:} \emph{turattu}, once more
  \item \emph{Say it:} \emph{nāy pūṉaiyait turattiyatu.} — The dog chased the cat.
  \item \emph{Recall:} say the Tamil for to hurl, then the Tamil for a race, then say \emph{turattu} again
\end{itemize}

\begin{tcolorbox}[breakable,colback=teal!4,colframe=teal!35!black,title={Before you move on}]
What does \ta{துரத்து} mean? (\textquotedblleft{}To chase; to drive away\textquotedblright{}.) Say it once more.
\end{tcolorbox}
\section[\texorpdfstring{\ta{சறுக்கு} (caṟukku)}{சறுக்கு (caṟukku)}]{\ta{சறுக்கு} (caṟukku) — to slide; to skid}

\label{lesson:TA-C533-carukku}



\begin{quote}
Before the new one: say the Tamil for a race, then the Tamil for to chase.
\end{quote}

\subsection*{You'll want to know: \ta{சறுக்கு}}
\textbf{\ta{சறுக்கு}} — \emph{caṟukku} — \textquotedblleft{}to slide; to skid\textquotedblright{}.

\begin{quote}
\ta{குழந்தைகள்} \ta{பனியில்} \ta{சறுக்கி} \ta{விளையாடினார்கள்}.
\end{quote}

\begin{quote}
\emph{kuḻantaikaḷ paṉiyil caṟukki viḷaiyāṭiṉārkaḷ.}
\end{quote}

\begin{quote}
The children played sliding on the snow.
\end{quote}

Past (I …): \ta{சறுக்கினேன்}, \emph{caṟukkiṉēṉ}.

Close to \ta{வழுக்கு} (\emph{vaḻukku}), to slip; \ta{சறுக்கு} is often a slide on purpose.

\begin{cousinweb}
\textbf{\ta{சறுக்கு}} (\emph{caṟukku}) is native Dravidian: DEDR 2482 gives \ta{சறுக்கு} (\emph{caṟukku}), \textquotedblleft{}to slip or slide\textquotedblright{}.
\end{cousinweb}

\begin{grammarlens}[title={What you've built}]
One more word for describing what you lived and felt.
\end{grammarlens}

\subsection*{Your turn}
\begin{itemize}
\raggedright
  \item \emph{Say it:} \emph{caṟukku}
  \item \emph{Say it:} \emph{caṟukku}, once more
  \item \emph{Say it:} \emph{kuḻantaikaḷ paṉiyil caṟukki viḷaiyāṭiṉārkaḷ.} — The children played sliding on the snow.
  \item \emph{Recall:} say the Tamil for a race, then the Tamil for to chase, then say \emph{caṟukku} again
\end{itemize}

\begin{tcolorbox}[breakable,colback=teal!4,colframe=teal!35!black,title={Before you move on}]
What does \ta{சறுக்கு} mean? (\textquotedblleft{}To slide; to skid\textquotedblright{}.) Say it once more.
\end{tcolorbox}
\section[\texorpdfstring{\ta{தாவு} (tāvu)}{தாவு (tāvu)}]{\ta{தாவு} (tāvu) — to leap; to jump across}

\label{lesson:TA-C533-tavu}



\begin{quote}
Before the new one: say the Tamil for to chase, then the Tamil for to slide.
\end{quote}

\subsection*{You'll want to know: \ta{தாவு}}
\textbf{\ta{தாவு}} — \emph{tāvu} — \textquotedblleft{}to leap; to jump across\textquotedblright{}.

\begin{quote}
\ta{குரங்கு} \ta{ஒரு} \ta{மரத்திலிருந்து} \ta{இன்னொரு} \ta{மரத்துக்குத்} \ta{தாவியது}.
\end{quote}

\begin{quote}
\emph{kuraṅku oru marattiliruntu iṉṉoru marattukkut tāviyatu.}
\end{quote}

\begin{quote}
The monkey leapt from one tree to another.
\end{quote}

Past (I …): \ta{தாவினேன்}, \emph{tāviṉēṉ}.

Where it leaps to takes -\ta{க்கு}. \ta{குதி} (\emph{kuti}) is a jump up and down; \ta{தாவு} is a leap across.

\begin{cousinweb}
\textbf{\ta{தாவு}} (\emph{tāvu}) is native Dravidian: DEDR 3177 gives \ta{தாவு} (\emph{tāvu}), \textquotedblleft{}to jump up, to leap over\textquotedblright{}.
\end{cousinweb}

\begin{grammarlens}[title={What you've built}]
One more word for describing what you lived and felt.
\end{grammarlens}

\subsection*{Your turn}
\begin{itemize}
\raggedright
  \item \emph{Say it:} \emph{tāvu}
  \item \emph{Say it:} \emph{tāvu}, once more
  \item \emph{Say it:} \emph{kuraṅku oru marattiliruntu iṉṉoru marattukkut tāviyatu.} — The monkey leapt from one tree to another.
  \item \emph{Recall:} say the Tamil for to chase, then the Tamil for to slide, then say \emph{tāvu} again
\end{itemize}

\begin{tcolorbox}[breakable,colback=teal!4,colframe=teal!35!black,title={Before you move on}]
What does \ta{தாவு} mean? (\textquotedblleft{}To leap; to jump across\textquotedblright{}.) Say it once more.
\end{tcolorbox}
\section[\texorpdfstring{\ta{முக்குளி} (mukkuḷi)}{முக்குளி (mukkuḷi)}]{\ta{முக்குளி} (mukkuḷi) — to dive}

\label{lesson:TA-C533-mukkuli}



\begin{quote}
Before the new one: say the Tamil for to slide, then the Tamil for to leap.
\end{quote}

\subsection*{You'll want to know: \ta{முக்குளி}}
\textbf{\ta{முக்குளி}} — \emph{mukkuḷi} — \textquotedblleft{}to dive\textquotedblright{}.

\begin{quote}
\ta{அவன்} \ta{குளத்தில்} \ta{முக்குளித்து} \ta{ஒரு} \ta{காசை} \ta{எடுத்தான்}.
\end{quote}

\begin{quote}
\emph{avaṉ kuḷattil mukkuḷittu oru kācai eṭuttāṉ.}
\end{quote}

\begin{quote}
He dived into the pond and picked up a coin.
\end{quote}

Past (I …): \ta{முக்குளித்தேன்}, \emph{mukkuḷittēṉ}.

Where one dives takes -\ta{இல்}: \ta{குளத்தில்} \ta{முக்குளி}.

\begin{cousinweb}
\textbf{\ta{முக்குளி}} (\emph{mukkuḷi}) ends in the native \ta{குளி} (\emph{kuḷi}), \textquotedblleft{}to bathe, to dive for pearls\textquotedblright{} (DEDR 1832); the origin of the first part, \ta{முக்கு}, is not settled.
\end{cousinweb}

\begin{grammarlens}[title={What you've built}]
One more word for describing what you lived and felt.
\end{grammarlens}

\subsection*{Your turn}
\begin{itemize}
\raggedright
  \item \emph{Say it:} \emph{mukkuḷi}
  \item \emph{Say it:} \emph{mukkuḷi}, once more
  \item \emph{Say it:} \emph{avaṉ kuḷattil mukkuḷittu oru kācai eṭuttāṉ.} — He dived into the pond and picked up a coin.
  \item \emph{Recall:} say the Tamil for to slide, then the Tamil for to leap, then say \emph{mukkuḷi} again
\end{itemize}

\begin{tcolorbox}[breakable,colback=teal!4,colframe=teal!35!black,title={Before you move on}]
What does \ta{முக்குளி} mean? (\textquotedblleft{}To dive\textquotedblright{}.) Say it once more.
\end{tcolorbox}
